\documentclass[12pt]{article}
\usepackage[utf8]{inputenc}
\usepackage[T1]{fontenc}
\usepackage{hyperref}
\title{Ternate Language Resources for K-12 Curriculum Development}
\author{Compiled by Claude Sonnet 5 (claude-sonnet-5) for Andr\'as Kornai}
\date{}
\begin{document}
\maketitle

\section{Introduction}
Ternate (ISO 639-3 tft; Glottocode tern1247) is a North Halmahera language of Maluku, Indonesia, with roughly 42,000 speakers. Searches found little material usable in schools. Much of what turns up concerns Ternate Malay, a different language (a Malay-based creole), and must not be confused with Ternate.

\section{Primary Educational Resources}
No primers, textbooks or school programmes for Ternate were found in the search results.

\section{Linguistic and Reference Materials}
\begin{itemize}
\item A Wiktionary contributor page lists a \emph{Kamus Ternate-Indonesia} (2009, Rusli Andi Atjo), \emph{Struktur Bahasa Ternate}, and Fortgens' Ternate texts. The search result gave only titles, with no publisher or availability details, so these need verification: \url{https://en.wiktionary.org/wiki/User:Alexlin01}
\item Journal article from Universitas Khairun on Ternate affixation (prefixes, infixes, suffixes, confixes): \url{https://ejournal.unkhair.ac.id/index.php/bilingual/article/view/4639/2974}
\item Verbix overview of Ternate verb subject marking: \url{https://docs.verbix.com/Languages/Ternate}
\item Leiden University repository items surfaced for this query. Their content was not confirmed to be about Ternate, so verify before use: \url{https://scholarlypublications.universiteitleiden.nl/access/item%3A2930091/view}
\item PARADISEC field notes page (North Halmahera materials, contents not verified): \url{https://www.paradisec.org.au/fieldnotes/INHALM.htm}
\item Related but different language: \emph{Ternate Malay: Grammar and Texts} (Ternate Malay, not Ternate): \url{https://lingweb.eva.mpg.de/archive/jakarta/ternate.php.html}
\end{itemize}

\section{Religious and Community Resources}
None found in the search results.

\section{Digital and Audio Resources}
No audio, video or software for Ternate were confirmed. The PARADISEC page above may point to archived recordings, but this was not verified.

\section{Curriculum Development Context}
The language is reported as endangered and not known to be taught in schools. The Khairun article notes that Ternate Malay has replaced much of its vocabulary. A curriculum would have to start from linguistic description and the Indonesian-language dictionary and grammar titles, and would need community and academic partners. No standard school orthography was identified.
\end{document}